
qui
sont des isomorphismes aux points @@M0@@. Considérons alors le complexe:
@@M1@@
La classe de @@M2@@ est @@M3@@; par ailleurs, @@M4@@ est acyclique aux points @@M5@@, donc @@M6@@ ne contient aucune composante irréductible de @@M7@@, et par conséquent est de codimension @@M8@@ dans @@M9@@; ceci montre que @@M10@@ est dans @@M11@@.

Soit maintenant @@M12@@ un élément de @@M13@@; nous allons montrer que @@M14@@. On peut supposer que @@M15@@ s'écrit:
@@M16@@
les @@M17@@, étant les classes de faisceaux localement libres de type fini, et, pour tout @@M18@@, @@M19@@ et @@M20@@ ayant même rang. De plus, quitte à décomposer @@M21@@ en somme d'un nombre fini d'ouverts, et à raisonner dans chaque ouvert, on peut supposer les @@M22@@ de rang constant. Soit alors @@M23@@ le produit fibré sur @@M24@@ des schémas de drapeaux des @@M25@@. Dans @@M26@@, les images des @@M27@@ sont sommes de classes de faisceaux inversibles. Soient @@M28@@ les classes de deux @@M29@@-Modules inversibles; on a:
@@M30@@
Donc @@M31@@. Or @@M32@@ est noethérien et a un faisceau ample, donc, d'après ce qui précède et 6.6, @@M33@@. Utilisant l'additivité de @@M34@@ et 6.6, on en déduit que @@M35@@, et @@M36@@, @@M37@@ étant le morphisme canonique @@M38@@.

Nous allons en déduire que @@M39@@. Soit @@M40@@ la dimension relative de @@M41@@ sur @@M42@@. Il existe un élément @@M43@@ de @@M44@@, appartenant à une base de @@M45@@ sur @@M46@@, tel que:
